% =============================================================================
% SEÇÃO 3: METODOLOGIA
% Diretrizes do Professor (Seções 3.2 e 3.3 do edital):
% - A "receita de bolo" técnica do seu projeto.
% - Origem e natureza dos dados utilizados (dataset, formato, pré-processamento).
% - Rigoroso processo de modelagem: como os elementos do problema real foram
%   traduzidos computacionalmente em vértices (V) e arestas (E) e/ou nós de uma árvore?
% - Descrição da lógica do algoritmo proposto passo a passo.
% - OBRIGATÓRIA A INCLUSÃO DE UMA FIGURA CONTENDO O FLUXOGRAMA DA SUA PROPOSTA,
%   mapeando o ciclo de vida da solução da entrada de dados até a obtenção dos resultados.
% =============================================================================

\section{Metodologia}
\label{sec:metodologia}

% TODO: Escreva a Metodologia aqui:
% 1. Origem e natureza da base de dados;
% 2. Processo formal de modelagem matemática (Definição de Vértices, Arestas e Pesos);
% 3. Lógica algorítmica passo a passo;
% 4. Inclusão e discussão do Fluxograma Obrigatório.


A metodologia deste estudo organiza-se em um pipeline estruturado em três eixos: (i) a ingestão, limpeza  e consolidação da base dados; (ii) modelagem relacional das vertentes de investigação e suas respcecitvas modelagens computacional baseadas em grafos; (iii) enquadramento e inferência dos resultados.


\subsection{Banco de Dados e Pipeline de Ingestão}

A obtenção de dados sobre compras públicas municipais no Brasil frequentemente esbarra na fragmentação dos portais locais. No município de Divinópolis/MG, o portal de transparência municipal~\cite{portal_transparencia_betha} fornece apenas consultas agregadas, impedindo o rastreamento atômico de lances e certames. Em contrapartida, o Portal de Dados Abertos do Tribunal de Contas do Estado de Minas Gerais (TCE-MG)~\cite{tce_mg_dados_abertos} centraliza as remessas formais obrigatórias do Sistema Informatizado de Contas dos Municípios (SICOM). Para este estudo, extraiu-se o conjunto completo de dados jurisdicionados brutos do município referentes ao exercício financeiro de 2025 integralmente consolidado.

Na sua forma bruta original, a base compreende 55 tabelas em formato CSV (delimitadas por ponto e vírgula), distribuídas em cinco módulos orçamentários: estrutura organizacional (\texttt{orgao}), certames e disputas (\texttt{licitacao}), instrumentos contratuais (\texttt{contrato}), reservas orçamentárias (\texttt{empenho}) e liquidações financeiras (\texttt{despesa}). A manipulação direta desses arquivos impunha alta redundância de metadados protocolares de envio e ruídos cadastrais (como perda de zeros à esquerda em documentos fiscais e valores monetários em texto regional).

Para viabilizar a modelagem computacional, desenvolveu-se um pipeline de Extração, Transformação e Carga (ETL) utilizando o motor analítico colunar \texttt{DuckDB} integrado a um banco relacional SQLite (\texttt{sicom.db}). A transformação guiou-se por três operações centrais: (i) projeção seletiva das entidades e tabelas relevantes para a rede (\texttt{secretarias}, \texttt{licitacoes}, \texttt{itens}, \texttt{propostas\_lances}, \texttt{homologacoes} e \texttt{empresas}); (ii) padronização vetorial via \texttt{LPAD} para recompor os 14 dígitos de CNPJs e 11 de CPFs, eliminando a criação espúria de vértices duplicados; e (iii) conversão de tipos de dados (datas em padrão ISO-8601 e valores monetários para ponto flutuante de dupla precisão \texttt{FLOAT64}). Esse saneamento reduziu em mais de 80\% o volume de dados redundantes, gerando uma base canônica e indexada para a construção dos grafos.

\subsection{Modelagem relacional e computacional dos grafos}


A investigação do ciclo de compras públicas municipais foi estruturada em três eixos analíticos complementares, cada um responsável por capturar uma dimensão específica do comportamento dos agentes econômicos e da administração pública. A partir do banco de dados canônico (sicom.db), cada eixo foi   formalmente traduzido em estruturas topológicas de grafos, definindo-se rigorosamente seus conjuntos de vértices, arestas, direcionamentos e funções de   ponderação

\subsubsection{}


% -----------------------------------------------------------------------------
% EXEMPLO DE INCLUSÃO DA FIGURA DO FLUXOGRAMA (OBRIGATÓRIO):
% -----------------------------------------------------------------------------
% \begin{figure}[!htbp]
% \centering
% \includegraphics[width=\columnwidth]{figures/fluxograma.png}
% \caption{Fluxograma metodológico do ciclo de vida da solução computacional proposta.}
% \label{fig:fluxograma}
% \end{figure}
